\chapter*{Kết luận và kiến nghị}
\addcontentsline{toc}{chapter}{Kết luận và kiến nghị}

\section*{1. Kết luận khoa học}
Đề tài nghiên cứu khoa học \textbf{"Nghiên cứu và phát triển khung mô hình thu nhận hình ảnh và dự báo dòng giao thông đô thị đa bước hai giai đoạn dựa trên mạng đồ thị không -- thời gian từ camera giám sát"} được thực hiện nhằm giải quyết một bài toán quan trọng trong quản lý đô thị thông minh: khai thác hiệu quả hạ tầng camera quan sát công cộng để thiết lập mạng lưới giám sát và dự báo lưu lượng giao thông tin cậy trong điều kiện dòng xe máy hỗn hợp tại các đô thị lớn của Việt Nam.

Qua quá trình nghiên cứu lý thuyết, mô hình hóa toán học và thực nghiệm trên quy mô toàn thành phố, đề tài đã hoàn thành các mục tiêu đặt ra với những kết quả cụ thể:

\begin{enumerate}
    \item \textbf{Xây dựng và công bố bộ dữ liệu chuẩn mực IC4SD-Traffic-HCM:} 
    Đề tài đã hoàn thành việc xây dựng bộ dữ liệu giao thông quy mô đô thị mở đầu tiên tại Việt Nam bao gồm \textbf{608 nút camera giám sát thực tế} tại TP. Hồ Chí Minh, bao gồm:
    \begin{itemize}
        \item 5.012 ảnh tĩnh độ phân giải cao được gán nhãn thủ công theo hai lớp ô tô con và xe máy, đạt hệ số tương quan nội lớp liên thẩm định viên cao ($\text{ICC} = 0.973, p < 0.001$).
        \item Chuỗi thời gian liên tục trong 86 ngày với 24.768 bước quan sát (chu kỳ 5 phút) bảo toàn phân loại phương tiện, được làm mượt bằng bộ lọc trượt kép nhân quả không gây rò rỉ dữ liệu tương lai.
        \item Đồ thị mạng lưới đường bộ OpenStreetMap với 2.452 cạnh định hướng liên thông thực tế, phản ánh chính xác cấu trúc hình học và quy định lưu thông của các tuyến đường.
    \end{itemize}
    Bộ dữ liệu và mã nguồn cài đặt được công bố công khai, tạo môi trường kiểm chuẩn chuẩn mực cho cộng đồng học thuật.

    \item \textbf{Hiệu quả của giải pháp hồi quy đếm xe trực tiếp (Giai đoạn 1):} 
    Nghiên cứu đã khẳng định tính hiệu quả của phương pháp hồi quy đếm xe trực tiếp dựa trên mạng xương sống EfficientNet-B5 \cite{efficientnet} kết hợp đầu hồi quy MLP thống nhất và hàm mất mát Huber Loss \cite{huber}. Cách tiếp cận này hạn chế ảnh hưởng của việc che khuất tương hỗ trong dòng xe máy và loại bỏ hoàn toàn chi phí tính toán của thuật toán NMS truyền thống, đạt sai số $\text{MAE} = 3.6363$ và tốc độ suy luận $35.2$ ms/batch, tiết kiệm $43.0\%$ lượng phép tính FLOPs so với ResNet-50 \cite{resnet}. Bản đồ nhiệt Grad-CAM++ xác nhận mô hình tập trung chính xác vào thân xe thay vì cảnh nền.

    \item \textbf{Đề xuất kiến trúc đồ thị cải tiến TA-STGCN (Giai đoạn 2):} 
    Đề tài đã phát triển mô hình tích hợp cơ chế chú ý tự thân theo thời gian đa đầu (Multi-Head Temporal Self-Attention) \cite{attention} xen kẽ giữa hai khối ST-Conv \cite{stgcn} trên đồ thị OpenStreetMap. Mô hình đạt sai số dự báo $\text{MAE} = 3.2401 \pm 0.0108$, duy trì tính ổn định qua 5 seed ngẫu nhiên ($\text{CV} = 0.33\%$) và kiểm soát tốt mức tăng sai số ở tầm nhìn dự báo xa 30 phút ($t+6$).

    \item \textbf{Đóng góp về lý thuyết lan truyền sai số và độ bền vững trước nhiễu:} 
    Đây là đóng góp khoa học trọng tâm của công trình. Đề tài đã chứng minh bằng giải tích toán học và kiểm chứng thực nghiệm rằng cơ chế chú ý theo trục thời gian hoạt động như một bộ lọc thông thấp thích ứng dữ liệu, làm giảm phương sai nhiễu tỷ lệ nghịch với độ rộng cửa sổ chú ý hiệu dụng ($\text{Var} = \sigma_\eta^2 / T_{\text{eff}}$). Dưới tác động của sai số nhận dạng mô phỏng thực tế (phân phối Gamma dị biệt có độ lệch âm do che khuất xe máy), TA-STGCN duy trì tốc độ suy thoái sai số thấp nhất, vượt trội rõ rệt các mô hình không có cơ chế chú ý thời gian (chỉ số Cohen's $d$ lên tới $+5.84$ so với Graph WaveNet \cite{wavenet} ở mức nhiễu cực hạn) và đạt độ bền bỉ tương đương các mô hình Transformer lớn (như STAEformer \cite{staeformer}) trong khi tiết kiệm $25.2\%$ bộ nhớ GPU VRAM và $70.4\%$ số lượng tham số so với ASTGCN \cite{astgcn}.

    \item \textbf{Tính khả thi của kiến trúc triển khai Edge--Cloud:} 
    Đề tài đã xây dựng phương án triển khai phân tán giúp tiết kiệm băng thông truyền thông tới \textbf{$62.500$ lần} (từ 437 GB/ngày xuống còn 7 MB/ngày) và bảo vệ quyền riêng tư hình ảnh theo nguyên tắc Privacy-by-Design, đồng thời ước tính tiết kiệm $89.5\%$ tổng chi phí sở hữu trong 5 năm (TCO) so với giải pháp vòng từ cảm ứng.
\end{enumerate}

\section*{2. Kiến nghị và đề xuất ứng dụng thực tiễn}
Từ các kết quả nghiên cứu và bài học thực tiễn thu được, nhóm tác giả đề xuất một số kiến nghị nhằm ứng dụng kết quả vào thực tế:

\subsection*{2.1. Kiến nghị đối với cơ quan quản lý giao thông và chính quyền đô thị}
\begin{enumerate}
    \item \textbf{Tận dụng hạ tầng camera sẵn có thay vì đầu tư mới vòng từ:} 
    Đề xuất các cơ quan quản lý giao thông tại các đô thị lớn xem xét áp dụng khung mô hình hai giai đoạn để tận dụng hệ thống camera giám sát hiện hữu, chuyển đổi thành mạng lưới cảm biến đếm xe và dự báo lưu lượng thông minh với chi phí đầu tư hạ tầng thấp hơn nhiều so với việc xẻ đường chôn vòng từ.
    \item \textbf{Thử nghiệm điều khiển đèn tín hiệu thích ứng đón đầu:} 
    Tích hợp luồng dữ liệu dự báo 15--30 phút của mô hình TA-STGCN vào các hệ thống điều khiển đèn tín hiệu tại các giao lộ trọng điểm để tự động điều chỉnh chu kỳ đèn đón đầu dòng xe, giảm thiểu nguy cơ hình thành và lan truyền sóng xung kích ùn tắc.
    \item \textbf{Áp dụng mô hình tính toán biên bảo vệ quyền riêng tư:} 
    Triển khai tính toán biên tại camera để giảm chi phí thuê băng thông truyền dẫn hàng tháng, đồng thời bảo đảm tuân thủ các quy định về bảo vệ dữ liệu cá nhân theo Nghị định 13/2023/NĐ-CP.
\end{enumerate}

\subsection*{2.2. Định hướng cho các nghiên cứu tiếp nối}
\begin{enumerate}
    \item \textbf{Mở rộng dữ liệu bằng học bán giám sát và mô hình nền tảng:} 
    Khuyến khích các nhóm nghiên cứu tiếp nối ứng dụng kỹ thuật học bán giám sát hoặc các mô hình nền tảng thị giác (Vision Foundation Models) để tự động sinh giả nhãn và mở rộng quy mô bộ dữ liệu ảnh.
    \item \textbf{Tích hợp đồ thị động thích ứng và dữ liệu đa phương thức:} 
    Kết hợp thêm các yếu tố ngoại sinh như lượng mưa thời gian thực, triều cường gây ngập cục bộ và lịch sự kiện đô thị vào đồ thị dự báo để nâng cao độ chính xác trong các tình huống thời tiết bất thường.
    \item \textbf{Đo kiểm trên phần cứng nhúng thực tế:} 
    Tiến hành tối ưu hóa mô hình sang định dạng TensorRT/INT8 và đo kiểm thực tế về tốc độ khung hình, độ trễ và công suất tiêu thụ trên các bo mạch biên chuyên dụng như NVIDIA Jetson Orin Nano, Raspberry Pi 5 và Google Coral TPU.
\end{enumerate}
